
\section{Introduction}

\subsection{Motivation}

The quadratic cost of transformer self-attention -- $\mathcal{O}(L^2)$ in 
both time and memory -- has motivated a decade of research into efficient 
alternatives. Among these, \textbf{linear recurrent attention} is particularly 
attractive because it reduces complexity to $\mathcal{O}(L)$ while 
maintaining the same asymptotic expressiveness.

The core idea: replace the softmax attention matrix with a **fixed-size 
state** $S \in \mathbb{R}^{d_h \times d_h}$ that is updated as tokens arrive:

\[
S_t = f(S_{t-1}, K_t, V_t), \quad y_t = Q_t^\top S_t
\tag{1.1}
\]

For this to work, $f$ must be \textbf{associative} -- that is, the composition 
$(f \circ f)$ must be efficiently computable.

0.85


\subsection{Problem Statement}

We seek an attention mechanism $\mathcal{A}$ with the following properties:

\begin{enumerate}
\item \textbf{Linear complexity:} $\mathcal{O}(L)$ time and memory.
\item \textbf{Bidirectional support:} work on encoders (BERT, RoBERTa).
\item \textbf{Causal support:} work on decoders (GPT-2, LLaMA).
\item \textbf{Weight compatibility:} transfer pretrained Q/K/V/O projections.
\item \textbf{Gating:} input-dependent forgetting/retention.
\item \textbf{Fine-tunable:} LoRA recovery.

\end{enumerate}
\subsection{Our Approach}

\textbf{Gated DeltaNet} combines three ideas:

\begin{enumerate}
\item \textbf{Delta Rule update} (Widrow--Hoff, 1960): correct the state by the error
  between the predicted value and the observed value:
  $$
  S\_t = S\_{t-1} + K\_t^\top (V\_t - K\_t S\_{t-1})
  \tag{1.2}
  $$

\item \textbf{Input-dependent gating:} replace fixed step size with learnable
  $\alpha_t, \beta_t \in (0, 1)$:
  $$
  S\_t = \beta\_t S\_{t-1} + \alpha\_t K\_t^\top (V\_t - K\_t S\_{t-1})
  \tag{1.3}
  $$

\item \textbf{Bidirectional scan for encoders:} average forward and backward scans:
  $$
  y = \frac{1}{2}(y\_{\text{fwd}} + y\_{\text{bwd}})
  \tag{1.4}
  $$

\end{enumerate}
\subsection{Contributions}

\begin{itemize}
\item \textbf{A gated delta rule attention} with input-dependent decay and
  bidirectional support.
\item \textbf{A chunked implementation} that is exact (not approximate) but uses
  nested Python loops.
\item \textbf{An honest evaluation} reporting that the implementation is
  \textbf{~45$\times$ slower} than baseline on DistilBERT and exceeds memory on larger
  models.
\item \textbf{A negative result:} the algorithm is not the bottleneck; the
  implementation is.
\item \textbf{Full reproducibility}, including code, timing measurements, and
  recommendations for fixing.

\end{itemize}
\subsection{Paper Organization}

Section 2 reviews linear attention, state-space models, and DeltaNet. 
Section 3 provides background on RNNs, associative scans, and the Delta 
Rule. Section 4 presents Gated DeltaNet. Section 5 covers architecture 
integration. Section 6 details the experimental setup. Section 7 reports 
results (including timeouts). Section 8 analyzes the implementation 
bottleneck. Section 9 discusses practical guidance. Section 10 lists 
limitations. Section 11 concludes.



\subsection{Notation}

\begin{table}[H]
\centering
\small
\begin{tabular}{lll}
\toprule
\textbf{Symbol} & \textbf{Meaning} & \textbf{First appears} \\
\midrule
$L$ & Sequence length & 1.1 \\
$d$ & Head dimension & 1.1 \\
$H$ & Number of heads & 1.1 \\
$S_t$ & State matrix at time $t$ & 1.1 \\
$\alpha_t$ & Forget gate (input-dependent) & 1.3 \\
$\beta_t$ & Write gate (input-dependent) & 1.3 \\
$K_t, V_t, Q_t$ & Key, Value, Query at token $t$ & 1.1 \\
$r$ & LoRA rank & 3.5 \\
$\text{chunk\_size}$ & Inner loop size & 4.4 \\
\bottomrule
\end{tabular}
\end{table}
</parameter>

